% SOURCE_STATEMENT: "Is the monodromy of the trigonometric Casimir connection equivalent to the quantum Weyl-group action of the affine braid group on finite-dimensional modules over the quantum loop algebra $U_{\\hbar}(L\\mathfrak g)$?"
% Preparation and final self-audit: CONTRIBUTING.md prompts C and D; source/audit artifacts retained in the local physics_statements_2026_10_11 batch.
% !TeX program = lualatex
\documentclass[11pt]{article}
\usepackage[margin=25mm]{geometry}
\usepackage{fontspec}
\setmainfont{Latin Modern Roman}
\usepackage{amsmath,amssymb,mathtools}
\usepackage{unicode-math}
\setmathfont{Latin Modern Math}
\usepackage{xurl,hyperref}
\hypersetup{colorlinks=true,urlcolor=blue}
\setcounter{secnumdepth}{0}
\setlength{\parindent}{0pt}
\setlength{\parskip}{5pt}
\setlength{\emergencystretch}{3em}
\newcommand{\sref}[2]{\hyperlink{source:#1:#2}{[S#2]}}
\newcommand{\cataloguescope}[1]{\paragraph{Recorded target.}#1}
\newcommand{\R}{\mathbb R}
\newcommand{\C}{\mathbb C}
\newcommand{\N}{\mathbb N}
\newcommand{\Z}{\mathbb Z}
\newcommand{\Q}{\mathbb Q}
\newcommand{\F}{\mathbb F}
\newcommand{\eps}{\varepsilon}
\begin{document}
\section{Casimir-Connection Monodromy and Quantum Weyl Groups}\label{statement}

\subsection{Definitions and mathematical statement}

\paragraph{Shared notation.}$\mathbb N=\{1,2,\ldots\}$; the standard number fields have their usual meanings. All additional conventions are stated locally.


\paragraph{Definitions and conventions.} Let $\mathfrak g$ be a complex simple Lie algebra, $\mathfrak h$ a Cartan subalgebra, and $W$ its Weyl group. The regular torus removes the root hypertori from the torus with Lie algebra $\mathfrak h$; the fundamental group of its Weyl quotient is the affine braid group in the report. The trigonometric Casimir connection is the flat Weyl-equivariant connection defined there with values in the Yangian $Y(\mathfrak g)$. Its monodromy gives a braid-group representation. The quantum loop algebra $U_\hbar(L\mathfrak g)$ deforms the enveloping algebra of $\mathfrak g[z,z^{-1}]$, and its quantum Weyl operators give a second braid-group action. Equivalence asks for an intertwiner between these representations under the source's module comparison.

\paragraph{Source scope and qualifications.} The identification of the finite-dimensional Yangian and quantum-loop module data is part of the proposed comparison. Equality of abstract braid groups alone is insufficient, and a result for sl\_2 alone does not settle every simple Lie algebra.

\paragraph{Catalogue mathematical target.}
\paragraph{Catalogue wording retained for traceability.}
Is the monodromy of the trigonometric Casimir connection equivalent to the quantum Weyl-group action of the affine braid group on finite-dimensional modules over the quantum loop algebra $U_{\hbar}(L\mathfrak g)$?

\paragraph{Saved literature context.} The supplied catalogue's assessment dated 2026-08-21 records: The monodromy equivalence is proved for sl\_2, while the conjecture for a general simple Lie algebra remains open. The assessment is retained as dated provenance; it is not a new comprehensive literature or open-status certification.
\paragraph{Scope and source qualification.}The identification of the finite-dimensional Yangian and quantum-loop module data is part of the proposed comparison. Equality of abstract braid groups alone is insufficient, and a result for sl_2 alone does not settle every simple Lie algebra.
\subsection{Short English statement}
Identify the monodromy of the trigonometric Casimir connection with the affine quantum Weyl-group action.
\subsection{Sources}
\begin{itemize}
\item[\hypertarget{source:statement:1}{S1}] ULAM.AI and the UnsolvedMath Contributors, UnsolvedMath, catalogue entry. Catalogue curation under CC BY 4.0; definitions and English exposition added locally.
\url{https://huggingface.co/datasets/ulamai/UnsolvedMath}.
\item[\hypertarget{source:statement:2}{S2}] Geometry, Quantum Fields, and Strings: Categorial Aspects (2010), Oberwolfach Reports, DOI 10.4171/owr/2010/25, pp. 1485--1533; official report source and mathematical context.
\url{https://doi.org/10.4171/owr/2010/25}.
\item[\hypertarget{source:statement:3}{S3}] Official EMS Press PDF of the same Oberwolfach report; PDF pages 29--30.
\url{https://ems.press/content/serial-article-files/46283}.
\item[\hypertarget{source:statement:4}{S4}] Toledano-Laredo; reference in the saved dated catalogue literature review.
\url{https://arxiv.org/abs/1003.2017}.
\item[\hypertarget{source:statement:5}{S5}] Gautam--Toledano-Laredo; reference in the saved dated catalogue literature review.
\url{https://arxiv.org/abs/1109.2367}.
\end{itemize}
\label{end:statement}
\end{document}
